\section*{Erd\H{o}s Problem 637: many distinct degrees in a Ramsey graph}

\paragraph{Statement and claim.}
An \(n\)-vertex graph with
\(\operatorname{hom}(G)\leq C\log n\) has an induced subgraph on
\(\Omega_C(n)\) vertices with \(\Omega_C(\sqrt n)\) distinct degrees.
Here \(\operatorname{hom}(G)\) is the maximum of the clique number
and independence number.
We reconstruct the diversity argument of Bukh and Sudakov,
\emph{Induced subgraphs of Ramsey graphs with many distinct degrees},
JCTB 97 (2007), 612--619:
\url{https://people.math.ethz.ch/~sudakovb/distinct-degrees.pdf}.
The final random sampling uses independent vertex choices.
See also \url{https://www.erdosproblems.com/637}.

\paragraph{External input.}
We use the fixed-density asymptotic form of the
Erd\H{o}s--Szemer\'edi homogeneous-set theorem (Theorem 2.1 of
that paper): for each fixed \(d_0\in(0,1/2]\), and all sufficiently
large \(M\) depending on \(d_0\), every graph \(H\) of order \(M\)
and edge density at most \(d_0\) satisfies
\[
 \operatorname{hom}(H)\geq
 c_0\,\frac{\log M}{d_0\log(1/d_0)}
\]
for an absolute positive constant \(c_0\).
Here \(d_0\) is fixed before \(M\) tends to infinity; this is not
a bound uniform down to arbitrarily small actual graph densities.
This is the only substantive external structural input.

\paragraph{A large diverse induced subgraph.}
Call an \(M\)-vertex graph \(\eta\)-diverse if each vertex has at most
\(M^{1/5}\) other vertices \(y\) for which
\(|N(x)\mathbin{\triangle}N(y)|<\eta M\).
Choose a sufficiently large integer \(K=K(C)\), and set
\(\eta=1/K\) and \(a=8^{-K}\).
We show that \(G\) contains an induced \(\eta\)-diverse graph of order
at least \(an\), for all sufficiently large \(n\).

Suppose otherwise. We recursively construct \(K\) disjoint sets
\(S_1,\ldots,S_K\), each of size
\(q=\lfloor(an)^{1/5}\rfloor\), and nested induced graphs
\(G=H_0\supset H_1\supset\cdots\supset H_K\), with
\(|H_i|\geq|H_{i-1}|/8\) and \(S_i\subset H_{i-1}\setminus H_i\).
In \(H=H_{i-1}\), failure of diversity supplies a vertex \(x\)
and \(q\) other vertices \(S_i\) whose neighbourhoods differ from
that of \(x\) in fewer than \(\eta|H|\) places.
In \(B=V(H)\setminus S_i\), which has size at least \(|H|/2\),
either the neighbours or the non-neighbours of \(x\) form a set
\(D\) of size at least \(|B|/2\).
Every vertex of \(S_i\) has the corresponding adjacency type to all
but at most \(\eta|H|\leq4\eta|D|\) vertices of \(D\).
By averaging, at least half the vertices of \(D\) have that type
to at least a \(1-8\eta\) fraction of \(S_i\).
Take these vertices for \(H_i\).

Thus every vertex of every later \(S_j\) has either at most
\(8\eta q\) neighbours in \(S_i\), or at least
\((1-8\eta)q\), with the choice depending only on \(i\).
Select at least \(K/2\) sets having the same choice, and complement
the graph if necessary. Their union \(S\) has density at most
\[
 8\eta+\frac{2}{K}=\frac{10}{K}.
\]
Also \(\log|S|\geq\tfrac15\log n-O_C(1)\).
Apply the external theorem with the fixed upper bound \(d_0=10/K\).
It gives
\[
 \operatorname{hom}(G)\geq
 \frac{c_0K}{10\log(K/10)}
       \bigl(\tfrac15\log n-O_C(1)\bigr)>C\log n
\]
if \(K\) was chosen sufficiently large and then \(n\) sufficiently
large. This contradiction proves the diversity assertion.

\paragraph{Sampling and counting collisions.}
Let \(H\) be the diverse graph, of order \(M\geq an\).
Choose a random subset \(W\) by retaining each vertex independently
with probability \(1/2\).
For a pair \(x,y\) with
\(|N_H(x)\triangle N_H(y)|\geq\eta M\), condition on retaining
both vertices. Their degree difference in \(H[W]\) is a signed sum
of \(D\geq\eta M-2\) independent fair Bernoulli variables.
Replacing each negatively signed variable \(X\) by \(1-X\)
shows that this sum is a translate of
\(\operatorname{Bin}(D,1/2)\).
Its largest atom is at most \(O(D^{-1/2})\), by the central
binomial coefficient estimate. Consequently this pair has equal
degrees with probability \(O_\eta(M^{-1/2})\).

There are at most \(M^{6/5}\) exceptional ordered pairs.
If \(Z\) counts unordered pairs of retained vertices with equal
degrees, it follows that
\[
 \mathbb EZ=O_\eta(M^{3/2}).
\]
Markov's inequality makes \(Z\leq4\mathbb EZ\) with probability
at least \(3/4\). Since \(|W|\) has mean \(M/2\) and variance
\(M/4\), Chebyshev's inequality gives
\(\Pr(|W|<M/3)\leq9/M\).
For sufficiently large \(M\), both desired events occur together.

For such a set \(W\), let the nonempty degree classes have sizes
\(u_1,\ldots,u_r\). Then
\[
 \sum_i u_i=|W|,\qquad
 \sum_i u_i^2=|W|+2Z=O_\eta(M^{3/2}).
\]
Cauchy--Schwarz gives
\[
 r\geq\frac{|W|^2}{|W|+2Z}
   =\Omega_\eta(\sqrt M)=\Omega_C(\sqrt n),
 \qquad |W|\geq an/3.
\]
This proves the requested assertion.

\paragraph{Scope.}
The proof reconstructs the known solution from the stated
Erd\H{o}s--Szemer\'edi theorem and elementary probabilistic estimates.
All constants depend only on \(C\); no uniform constant independent
of the Ramsey parameter is claimed.

\paragraph{Completion estimate.}
\textbf{COMPLETION: 100\%}
